\noindent
\textbf{Information-Seeking Support in NLP}\quad
NLP research supporting human information-seeking has mainly focused on building question-answering (QA) systems~\citep{chen-etal-2017-reading, lee-etal-2019-latent,dasigi-etal-2021-dataset,levy-etal-2021-open,yuan-etal-2020-interactive-machine}. These works often assume that the answer can be found within a single document~\citep{clark-etal-2020-tydi} or that users can formulate complex queries~\citep{yang-etal-2018-hotpotqa,chen2021open,ahmadvand2023making}, assumptions that do not hold true in complex information seeking~\citep{butler2000learners,booth2009craft,bystrom1995task}.

%However, in complex information seeking, prevalent in learning~\citep{butler2000learners}, research~\citep{booth2009craft}, and decision-making tasks~\citep{bystrom1995task}, relevant information usually come from multiple sources, and users may lack the knowledge to formulate appropriate queries or have evolving intents. 

Some more recent works have proposed long-form QA systems~\citep{xu-etal-2023-critical,xu2024kiwi} and automatic expository writing systems~\citep{balepur-etal-2023-expository,shen2023summarization,shao2024assisting} to synthesize information from multiple sources. %, generating outputs with good breadth and depth. 
Some other studies have explored conversational search~\citep{kumar-callan-2020-making,nakamura-etal-2022-hybridialogue}. %, which supports multi-turn interactions between users and the system. However, these works typically focus on seeking basic factual information or simply breaking down complex queries into multi-turn questions. 
However, these works typically ignore human interaction or only passively answer user questions. We construct a multi-agent system with a human-in-the-loop protocol to support effective user interaction for complex and evolving information needs.



% [used] Making Information Seeking Easier: An Improved Pipeline for Conversational Search (kumar-callan-2020-making): This paper presents a highly effective pipeline for passage retrieval in a conversational search setting. (Questions are straightforward. An example: Title: goat breeds; Description: Interested in buying goats that implies interest in different breeds of goats and their use; Turns: 1. What are the main breeds of goat? 2. Tell me about boer goats. 3. What breed is good for meat? 4: Are angora goats good for it? ...)

% [used] TYDI QA: ABenchmark for Information-Seeking Question Answering in Typologically Diverse Languages (clark-etal-2020-tydi): To be realistic and avoid priming effects, questions are written by people who want to know the answer, but don't know the answer yet. The expected answer is a span in Wikipedia article.

% [used] Interactive Machine Comprehension with Information Seeking Agents (yuan-etal-2020-interactive-machine): restrict the document context that a model observes at one time; the task requires models to ‘feed themselves’ rather than spoon-feeding them with information.

% [used] HybriDialogue: An Information-Seeking Dialogue Dataset Grounded on Tabular and Textual Data (nakamura-etal-2022-hybridialogue): crowdsourced natural conversations grounded on both Wikipedia text and tables; The conversations are created through the decomposition of complex multihop questions (e.g., What was the year of birth of the prime minister who had a Canadian state funeral in 1892 at Javis Street Baptist Church?) into simple, realistic multiturn dialogue interactions.

% [used] A Dataset of Information-Seeking Questions and Answers Anchored in Research Papers (dasigi-etal-2021-dataset): document grounded, information-seeking QA; Each question is written by an NLP practitioner who read only the title and abstract of the corresponding paper, and the question seeks information present in the full text.

% Stay Hungry, Stay Focused: Generating Informative and Specific Questions in Information-Seeking Conversations (qi-etal-2020-stay): 

% [used] KIWI: A Dataset of Knowledge-Intensive Writing Instructions for Answering Research Questions (xu2024kiwi): providing writing assistance to compose a long-form answer.



\noindent
\textbf{Multi-Agent Systems}\quad
As LMs advance, a growing body of research explores their use in multi-agent applications~\citep{wu2023autogen,BabyAGI,liu2023dynamic,wang2024unleashing}. Several studies show that multi-agent debate enhances factuality and reasoning compared to using a single LM~\citep{du2023improving,liang2023encouraging}, and cooperative role-playing frameworks improve performance on coding or mathematical benchmarks~\citep{NEURIPS2023_a3621ee9,hong2023metagpt}. While these studies primarily focus on automating tasks, the potential applications extend further. For instance, Generative Agents~\citep{Park2023GenerativeAgents} instantiate an interactive environment with twenty-five LM agents to study emergent social behaviors, and \citet{michael2023debate} show that multi-agent debates help humans supervise model outputs. Our work aligns with these broader applications by constructing a multi-agent system to facilitate human learning.
%Another line of work constructs multi-agent system through personifying LMs with different roles. \citet{NEURIPS2023_a3621ee9} proposes CAMEL, a cooperative role-playing communication framework, to improve completion rate on coding or math problems; MetaGPT~\citep{hong2023metagpt} proposes a multi-agent collaborative framework to improve software engineering benchmark; Generative Agents~\citep{Park2023GenerativeAgents} instantiates a sandbox environment with twenty-five LM agents with different characters to study emergent social behaviors. 

\noindent
\textbf{Collaborative Discourse for Human Learning}\quad
Collaborative discourse has long been valued in classroom settings for its ability to deepen learners' understanding of concepts, enhance peer learning, and increase engagement %in the learning process~
\citep{nussbaum2008collaborative,osborne2010arguing,kolodner2007roles,chinn2000structure}. Specifically, \citet{nussbaum2008collaborative} argues not all types of collaborative discourse are equally beneficial to students' learning, emphasizing the importance of critical discussion where participants assume different points of view. Furthermore, the facilitator role is important in collaborative discourse, with asking questions and providing complementary information as popular strategies~\citep{onrubia2022assisting}.
